\Needspace{8\baselineskip}
\Article{272}
We will treat the theorem proposed in the preceding article, that all ternary forms of a given determinant fall into a \emph{finite} number of classes, by a method similar to the one we used for binary forms. We will first show how any ternary form can be reduced to a simpler form, and then that for any given determinant there are only finitely many simplest forms (reached by such reductions). Suppose \SourceLeaf{313}in general that a ternary form \(f=\left(\begin{smallmatrix}a,&a',&a''\\b,&b',&b''\end{smallmatrix}\right)\) of determinant \(D\) (different from zero) passes, by the substitution \((S)\)
\[
\begin{array}{ccc}
\alpha,&\beta,&\gamma\\
\alpha',&\beta',&\gamma'\\
\alpha'',&\beta'',&\gamma''
\end{array}
\]
into an equivalent form \(g=\left(\begin{smallmatrix}m,&m',&m''\\n,&n',&n''\end{smallmatrix}\right)\). Our task will be to determine \(\alpha\), \(\beta\), \(\gamma\), etc.\ so that the form \(g\) becomes simpler than \(f\). Let the adjoints of \(f\), \(g\) be, respectively, \(\left(\begin{smallmatrix}A,&A',&A''\\B,&B',&B''\end{smallmatrix}\right)\), \(\left(\begin{smallmatrix}M,&M',&M''\\N,&N',&N''\end{smallmatrix}\right)\). denoted by \(F\), \(G\). Then by Article 269, \(F\) will pass into \(G\) by the adjoint of the substitution \(S\), and \(G\) into \(F\) by the substitution obtained by transposing \(S\). The number
\[
\alpha\beta'\gamma''+\alpha'\beta''\gamma+\alpha''\beta\gamma'-\alpha''\beta'\gamma-\alpha\beta''\gamma'-\alpha'\beta\gamma''
\]
which must be either \(=+1\) or \(=-1\), will be denoted by \(k\). With this understood, we observe:

I. If we set \(\gamma=0\), \(\gamma'=0\), \(\alpha''=0\), \(\beta''=0\), \(\gamma''=1\), we will have
\[
\begin{aligned}
m&=a\alpha\alpha+2b''\alpha\alpha'+a'\alpha'\alpha',\\
m'&=a\beta\beta+2b''\beta\beta'+a'\beta'\beta',\\
m''&=a''\\
n&=b\beta'+b'\beta,\\
n'&=b\alpha'+b'\alpha,\\
n''&=a\alpha\beta+b''(\alpha\beta'+\beta\alpha')+a'\alpha'\beta'
\end{aligned}
\]
Also, \(\alpha\beta'-\beta\alpha'\) must be either \(=+1\) or \(=-1\). It is clear from this that the binary form \((a,b'',a')\), whose determinant is \(A''\), is transformed by the substitution \(\alpha\), \(\beta\), \(\alpha'\), \(\beta'\) into the binary form \((m,n'',m')\) of determinant \(M''\), and is therefore equivalent to it, because \(\alpha\beta'-\beta\alpha'=\pm1\). Thus \(M''=A''\), which is also easily confirmed directly. So unless \((a,b'',a')\) is already a simplest form in its class, we may choose \(\alpha\), \(\beta\), \(\alpha'\), \(\beta'\) so that \((m,n'',m')\) is a simpler form. Indeed, the theory of equivalence of binary forms readily shows that this can be done so that \(m\) is no greater than \(\sqrt{-\frac{4}{3}A''}\) if \(A''\) is negative, or no greater than \(\sqrt{A''}\) if \(A''\) is positive, or \(m=0\) if \(A''=0\). Thus in every case the (absolute) value of \(m\) can certainly be lowered below, or at least as far as, \(\sqrt{\pm\frac{4}{3}A''}\). In this way the form \(f\) is reduced to another with a smaller first coefficient, if possible, and whose adjoint has the same third coefficient as the form \(F\) adjoint to \(f\). This is the \emph{first reduction}.

\SourceLeaf{314}
II. If, however, we set \(\alpha=1\), \(\beta=0\), \(\gamma=0\), \(\alpha'=0\), \(\alpha''=0\), we will have \(k=\beta'\gamma''-\beta''\gamma'=\pm1\). So the adjoint of the substitution \(S\) will be
\[
\begin{array}{rrr}
\pm1,&0,&0\\
0,&\gamma'',&-\beta''\\
0,&-\gamma',&\beta'
\end{array}
\]
by which \(F\) will pass into \(G\). We will therefore have
\[
\begin{aligned}
m&=a,\quad n'=b'\gamma''+b''\gamma',\quad n''=b'\beta''+b''\beta'\\
m'&=a'\beta'\beta'+2b\beta'\beta''+a''\beta''\beta''\\
m''&=a'\gamma'\gamma'+2b\gamma'\gamma''+a''\gamma''\gamma''\\
n&=a'\beta'\gamma'+b(\beta'\gamma''+\gamma'\beta'')+a''\beta''\gamma''\\
M'&=A'\gamma''\gamma''-2B\gamma'\gamma''+A''\gamma'\gamma'\\
N&=-A'\beta''\gamma''+B(\beta'\gamma''+\gamma'\beta'')-A''\beta'\gamma'\\
M''&=A'\beta''\beta''-2B\beta'\beta''+A''\beta'\beta'
\end{aligned}
\]
It follows that the binary form \((A'',B,A')\), whose determinant is \(Da\), passes by the substitution \(\beta'\), \(-\gamma'\), \(-\beta''\), \(\gamma''\) into the form \((M'',N,M')\) of determinant \(Dm\), and so is equivalent to it (because \(\beta'\gamma''-\gamma'\beta''=\pm1\), or because \(Da=Dm\)). So unless \((A'',B,A')\) is already a simplest form in its class, the coefficients \(\beta'\), \(\gamma'\), \(\beta''\), \(\gamma''\) can be chosen so that \((M''.\,N,M')\) is simpler. This can always be done so that \(M''\), without regard to sign, is no greater than \(\sqrt{\pm\frac{4}{3}Da}\). In this way the form \(f\) is reduced to another with the same first coefficient, but whose adjoint has, if possible, a smaller third coefficient than the form \(F\) adjoint to \(f\). This is the \emph{second reduction}.

III. Thus if \(f\) is a ternary form to which neither the first nor the second reduction is applicable, \emph{that is}, which cannot be transformed into a simpler form by either: we must have both \(aa<\) or \(=\frac{4}{3}A''\), and \(A''A''<\) or \(=\frac{4}{3}aD\), without regard to sign. Hence \(a^4\) will be \(<\) or \(=\frac{16}{9}A''A''\), and so \(a^4<\) or \(=\frac{64}{27}aD\), \(a^3<\) or \(=\frac{64}{27}D\), and \(a<\) or \(=\frac{4}{3}\sqrt[3]{D}\); from this again \(A''A''<\) or \(=\frac{16}{9}\sqrt[3]{D^4}\), and \(A''<\) or \(=\frac{4}{3}\sqrt[3]{D^2}\). So as long as \(a\) or \(A''\) still exceeds these bounds, one or the other of the preceding reductions must be applicable to the form \(f\). --- This conclusion must not be reversed, however, since it quite often
\SourceLeaf{315}
happens that a ternary form whose first coefficient, and the third coefficient of its adjoint, are already below those bounds can nevertheless be made still simpler by one or the other reduction.

\enlargethispage{-2\baselineskip}IV. If the first and second reductions are applied alternately to any given ternary form of determinant \(D\), \emph{that is}, the first or second to the form itself, the second or first to the resulting form, again the first or second to the next resulting form, etc., then clearly we must eventually reach a form to which neither can be applied any further. For the absolute magnitudes both of the first coefficients of the forms obtained in this way and of the third coefficients of their adjoints continually alternate between staying the same and decreasing. This process must eventually end somewhere, since otherwise there would be two infinite sequences of continually decreasing numbers. We have now obtained a remarkable theorem: \emph{Every ternary form of determinant \(D\) can be reduced to an equivalent form whose first coefficient is no greater than \(\frac{4}{3}\sqrt[3]{D}\), and the third coefficient of its adjoint no greater than \(\frac{4}{3}\sqrt[3]{D^2}\), without regard to sign, if the given form does not already have these properties.} --- Instead of the first coefficient of the form \(f\) and the third of the adjoint of \(f\), we could have treated in exactly the same way either the first coefficient of the form itself and the second of its adjoint; or the second of the form itself and the first or third of its adjoint; or the third of the form itself and the first or second of its adjoint. Each of these routes would equally have led us to our intended goal. But it is useful to keep consistently to one method, so that the operations involved can more easily be reduced to a fixed algorithm. Finally, we observe that for the two coefficients which we have shown how to lower below fixed bounds. still smaller bounds can be established if definite forms are separated from indefinite ones; this is not needed for our present purpose.
